\chapter{{\fontsize{14pt}{21pt}\selectfont\MakeUppercase{Приложение А. Техническое задание}}}
\label{app:tz}

\noindent{}Наименование программного изделия: программный комплекс корпоративного обучения с веб-панелью администрирования, мобильным клиентом слушателя и серверной частью REST API.

\noindent{}Обозначение: ПК «Coursum-LMS».\\
\noindent{}Основание для разработки: выпускная квалификационная работа по направлению 09.03.03 «Прикладная информатика», профиль «Корпоративные информационные системы».

\section*{А.1 Введение}

Наименование темы разработки: «Веб и мобильное приложение для корпоративного обучения с поддержкой адаптивного тестирования».

Объект автоматизации: процессы организации корпоративного обучения, назначения учебных программ, контроля освоения материала и оценки результатов в мультиарендной среде\cite{gost1920178,gost1930179,gost3460190}.

\section*{А.2 Назначение программы}

Программа предназначена для автоматизации функций корпоративной системы управления обучением: ведение пользователей и ролей в разрезе организаций-арендаторов; управление курсами, уроками и медиаконтентом; назначение обучения пользователям и группам; проведение адаптивного тестирования с изменением сложности заданий; формирование рекомендаций по повторению тем; предоставление аналитики по прогрессу и результатам.

\section*{А.3 Требования к программе}

\subsection*{А.3.1 Функциональные требования}

\begin{enumerate}
  \item система должна обеспечивать регистрацию, аутентификацию и выдачу токенов доступа и обновления;
  \item система должна поддерживать выбор текущей организации и проверку членства пользователя в организации;
  \item система должна реализовать роли слушателя, преподавателя, администратора организации и системного администратора с разграничением прав;
  \item система должна обеспечивать создание и редактирование курсов, уроков, тестов и вопросов с указанием сложности;
  \item система должна поддерживать назначение курсов пользователям и учебным группам;
  \item система должна обеспечивать запуск попытки тестирования, выдачу следующего вопроса с учётом текущей сложности, приём ответа и завершение попытки;
  \item система должна формировать рекомендации по темам для повторного изучения после завершения попытки;
  \item система должна предоставлять аналитические сводки по курсам и слушателям в веб-панели.
\end{enumerate}

\subsection*{А.3.2 Требования к надёжности}

\begin{enumerate}
  \item при недоступности СУБД серверная часть должна возвращать диагностируемый код ошибки;
  \item операции изменения данных должны выполняться в транзакциях с сохранением ссылочной целостности;
  \item журнал аудита должен фиксировать ключевые действия пользователей.
\end{enumerate}

\subsection*{А.3.3 Условия эксплуатации}

Клиентские приложения функционируют под управлением ОС семейств Windows, Linux (веб-браузер), Android/iOS (мобильное приложение). Серверная часть --- под управлением ОС Linux или Windows с поддержкой контейнеризации Docker.

\subsection*{А.3.4 Состав и параметры технических средств}

Минимальная конфигурация сервера: 2 ядра CPU, 4 Гбайт ОЗУ, 20 Гбайт диска; СУБД PostgreSQL не ниже версии 15 в составе контейнера.

\subsection*{А.3.5 Информационная совместимость}

Обмен данными между клиентами и серверной частью --- по протоколу HTTPS, формат тел сообщений JSON, базовый префикс API /api/v1.

\subsection*{А.3.6 Программная совместимость}

Серверная часть: интерпретатор Python 3, фреймворк FastAPI, ORM SQLAlchemy, миграции Alembic. Веб-клиент: среда выполнения JavaScript, библиотека React, язык TypeScript. Мобильный клиент: Flutter/Dart.

\subsection*{А.3.7 Маркировка и упаковка}

Поставка осуществляется в виде архива исходного кода или образов контейнеров с сопроводительным файлом версии и журналом изменений.

\subsection*{А.3.8 Требования по безопасности}

Хранение паролей --- только в виде криптографического хэша; доступ к данным арендатора --- только после проверки идентификатора организации и роли; использование TLS на границе контура эксплуатации.

\subsection*{А.3.9 Эксплуатационная документация}

Комплект включает руководство пользователя мобильного приложения, руководство администратора веб-панели, описание развёртывания (runbook), программу и методику испытаний.

\section*{А.4 Технико-экономические показатели}

Ожидаемый эффект внедрения связан со сокращением ручных операций при назначении обучения и учёте результатов по сравнению с разрозненными средствами; количественная оценка выходит за рамки настоящей работы и может быть выполнена отдельно по экономическим методикам организации-заказчика при пилотной эксплуатации.

\section*{А.5 Стадии и этапы разработки}

\begin{enumerate}
  \item техническое задание и уточнение требований;
  \item проектирование архитектуры и модели данных;
  \item реализация серверной части и REST API;
  \item реализация веб-панели и мобильного клиента;
  \item испытания и подготовка документации.
\end{enumerate}

\section*{А.6 Порядок контроля и приёмки}

Приёмка включает проверку соответствия функциональным требованиям по сценариям приложения~Б, проверку развёртывания по инструкции приложения~Г, анализ результатов автоматизированных тестов и контрольный просмотр журналов ошибок при демонстрации критического пути пользователя.
